\chapter{质量保证、测试体系与验收逻辑}

\section{测试层次}

项目测试约 3400 行，分为三层：

\begin{itemize}
  \item \term{单元测试}：配置解析、几何不变量、导管估计、锚点射线、Y 梁汇合点、特殊
  拓扑和文档 docstring；
  \item \term{Blender 后端测试}：封闭圆管扫掠、导管重建、体素并集、布尔差集、连通体
  清理和重复三角形检测；
  \item \term{端到端测试}：从病例配置运行生成流程，检查 STL 导出和验证结果。
\end{itemize}

单元测试特别覆盖了几个容易回归的几何决策：实心 cutter 不能被选为导管；上下 Q/P 的边缘
余量和偏移公式；U/背 U 射线角度；缺牙插值；低梁局部下潜；Y 梁 maximin 锚点；末端节点
共享；input 导管模式和 YAML 合并规则。

\section{生成阶段的失败关闭}

下列情况不会静默退化：不足两根有孔导管、导板标架退化、牙位报告 QA 不完整、上下颌或
输入路径不一致、射线找不到外壁出口、两导管内外侧差异不足、Y 梁三锚点近共线、布尔结果
不是封闭体、生产 review 状态未确认等。失败关闭能阻止几何错误进入最终 STL，但也意味着
新增病例通常需要明确修正语义和参数，不能只放宽容差。

\section{最终 STL 验证项目}

\begin{longtable}{p{0.27\textwidth}p{0.63\textwidth}}
  \toprule
  验证项 & 方法与阈值\\
  \midrule
  \endfirsthead
  \toprule
  验证项 & 方法与阈值\\
  \midrule
  \endhead
  拓扑 & 边界边、非流形边和重复三角形均为 0；恰有一个连通体。\\
  导管保留 & 参数化或输入导管表面样本位于最终体内，或距离表面不超过两个体素；保留率至少 0.90。\\
  主连接梁 & 排除孔腔内不应保留的中心线点后，中心线保留率至少 0.95；P 点和 A 点距离受限。\\
  根部强化 & 唯一导板端数量正确，球头探针和贴合脚探针具有足够包覆率。\\
  Y 梁 & 三臂中心线保留率、汇合点和应有的导板端强化。\\
  导孔 & 沿轴向自适应截面，在 0、0.5、0.9 倍可用半径上布置多角度探针；阻塞数必须为 0。\\
  观察窗 & 每条最终修正语义轴的点不在实体内，且到最近表面距离不小于 0.20 mm。\\
  特殊结构 & 根据模式检查跨分量、末端 U 梁或远中公共节点的几何公式和保留率。\\
  \bottomrule
\end{longtable}

\section{当前验证尚未覆盖的部分}

\begin{itemize}
  \item 手机包络在生成时被切除，但 \texttt{validate\_guide()} 不重新构造包络并逐点验证净空；
  \item 操作窗没有像导孔和观察窗一样形成独立、详细的最终探针报告；
  \item input 模式以配置内径作为孔道参考，不能替代对真实输入导管全孔腔的高精度计量；
  \item 结构检查主要是中心线和局部探针保留率，不直接计算有限元强度、疲劳、打印支撑或变形；
  \item 封闭和单连通不等同于没有自相交、过薄壁或不可清洁死角。
\end{itemize}

\section{本次报告构建期间的验证状态}

本次已读取现有过程报告、检查 tooth-47 图片，并用独立 Trimesh 读取最终 STL，确认其为单一
封闭有效体。尝试通过系统 Python 运行单元测试时缺少 pytest；随后尝试调用 Blender 5.2
后台测试环境，但 Blender 在启动后崩溃并写出 crash 文件，未得到可信的测试通过结果。
因此本报告不会把“测试文件存在”表述为“当前环境测试已全部通过”。

建议交接验收在稳定 Blender 环境中依次执行：文档严格构建、全部单元测试、Blender 后端
测试、tooth-47 当前版本重新生成，以及 \texttt{generate --validate}。每次运行都应保存
命令、退出码和验证 JSON。

